\documentclass{homework}
% \usepackage{graphicx} % Required for inserting images
\usepackage{amsmath}
\usepackage{gensymb}
\usepackage{subfig}
\usepackage{float}

\usepackage{listings}
\usepackage{color}


\class{APPM 5560}
\title{Homework 1}
\author{Sean Jennings}
\date{January 18, 2024}

\definecolor{dkgreen}{rgb}{0,0.6,0}
\definecolor{gray}{rgb}{0.5,0.5,0.5}
\definecolor{mauve}{rgb}{0.58,0,0.82}

\lstset{frame=tb,
  language=Python,
  aboveskip=3mm,
  belowskip=3mm,
  showstringspaces=false,
  columns=flexible,
  basicstyle={\small\ttfamily},
  numbers=none,
  numberstyle=\tiny\color{gray},
  keywordstyle=\color{blue},
  commentstyle=\color{dkgreen},
  stringstyle=\color{mauve},
  breaklines=true,
  breakatwhitespace=true,
  tabsize=3
}



\begin{document}
\maketitle

\section{Problem 1}
Each event $A$, $B$, and $C$ have probability
\begin{equation*}
    {\bf P}(A) = {\bf P}(B) = {\bf P}(C) = \frac{1}{2}
\end{equation*}
since they each correspond to 2 out of 4 possible outcomes. The 
intersection of each pair of events corresponds to the outcome
that Ball 1 is selected, and thus has probability
\begin{equation*}
    {\bf P}(A \cap B) = {\bf P}(A \cap C) = {\bf P}(B \cap C) = \frac{1}{4}.
\end{equation*}
Therefore, the conditions for independence
${\bf P}(A \cap B) = {\bf P}(A) {\bf P}(B)$, etc., for each pair of 
events. On the other hand, since the selection of Ball 1 implies the
occurence of each event $A$, $B$, and $C$, we also have
\begin{equation*}
    {\bf P}(A \cap B \cap C) = 0.25 \neq {\bf P}(A) {\bf P}(B) {\bf P}(C)
\end{equation*}
and thus the events are not altogether independent.

\section{Problem 2}

For $f(x)$ to be a probability distribution, it must satisfy the normalization
property. Specifically, the integral
of $f$ over all possible values of $x$ equal to 1:
\begin{align*}
    1= \int_{-\infty}^{\infty} f(x)dx &= \int_{1}^{\infty} Cx^{-\gamma} dx \\
        &= \frac{C}{-\gamma + 1} x^{-\gamma + 1} \eval_{x=1}^{\infty} \\
        &= \frac{C}{-\gamma + 1} (0 - 1) \\
        &= \frac{C}{\gamma - 1}.
\end{align*}
The value of $C$ for which this holds is $C = \gamma - 1$.

\section{Problem 3}

Let $p_h=0.65$ and $p_s=0.95$ be the probability of a hard and soft serve going in, respectively.
Let $q_h=0.7$ and $q_s=a$ be the probability of winning the point given the serve went in,
for a hard and soft serve, resp. Then, the conditional probability of winning the
point, given the hard serve chosen, is:
\begin{equation*}
    {\bf P}(\text{winning } | \text{ hard second serve}) = {\bf P}(\text{winning } | \text{ hard serve went in}) {\bf P}(\text{hard serve went it}) = q_h p_h.
\end{equation*}
Similarly, given a soft serve is chosen, the condition probability of winning is
\begin{equation*}
    {\bf P}(\text{winning } | \text{ soft second serve}) = q_s p_s = a p_s.
\end{equation*}
Thus, the probability of winning with the hard serve is greater than with the soft serve when
\begin{equation*}
    a < \frac{q_h p_h}{p_s} = \frac{0.7(0.65)}{0.95} \approx 0.479
\end{equation*}

\section{Problem 4}
\begin{letterenum}
\item Since $X$ is the number of coin flips until tails appears, $X$ is a geometric
R.V. We can find its MGF via recursion (given the first trial fails, the expectation is 
equivalent to the unconditional expectation, plus 1):
\begin{align*}
    \phi_X(s) &= {\bf E}[e^{sX}] \\
        &= pe^s + (1 - p) {\bf E}[e^{s(X+1)}] \\
        &= pe^s + (1 - p)e^s {\bf E}[e^{sX}].
\end{align*}
Grouping ${\bf E}[e^{sX}]$ terms yields:
\begin{equation*}
    (1 - e^s(1 - p)) {\bf E}[e^{sX}] = pe^s
\end{equation*}
so that the MGF is given by:
\begin{equation*}
    \phi_X(s) = {\bf E}[e^{sX}] = \frac{pe^s}{1 - e^s(1 - p)}
\end{equation*}

\item To find first and second moments of $X$, we use the quotient rule to take
derivatives of the MGF:
\begin{align*}
    \phi_X'(s) &= \frac{(1 - e^s(1-p))pe^s - pe^s (-e^s(1-p))}{(1 - e^s(1-p))^2} \\
        &= \frac{pe^s}{(1 - e^s(1-p))^2}
\end{align*}
Putting $s=0$, we have
\begin{equation*}
    {\bf E}[X] = \phi_X'(0) = \frac{p}{(1 - (1-p))^2} = \frac{1}{p}
\end{equation*}

The second derivative of the MGF is given by:
\begin{align*}
    \phi_X''(s) &= \frac{(1 - e^s(1-p))^2pe^s - (-e^s(1-p))2(1-e^s(1-p))pe^s}{(1-e^s(1-p))^4} \\
        &= \frac{pe^s(1 - e^s(1-p))\biggl(1 - e^s(1-p) + 2e^s(1-p)\biggr)}{(1 - e^s(1-p))^4} \\
        &= \frac{pe^s(1 + e^s(1-p))}{(1 - e^s(1-p))^3}
\end{align*}
Thus, the second moment is:
\begin{equation*}
    {\bf E}[X^2] = \phi_X''(0) = \frac{p(2-p)}{p^3} = \frac{2-p}{p^2}
\end{equation*}
\end{letterenum}

\section{Problem 5}

The PDF of a Poisson distribution with parameter $\lambda$ is given by
\begin{equation*}
    f_X(k) = e^{-\lambda} \frac{\lambda^k}{k!}.
\end{equation*}

Then, its MGF is
\begin{align*}
    \phi_X(s) &= {\bf E}[e^{sX}] \\
        &= \sum_{k=0}^\infty e^{sk} e^{-\lambda} \frac{\lambda^k}{k!} \\
        &= e^{-\lambda} \sum_{k=0}^\infty \frac{(\lambda e^s)^k}{k!} \\
        &= e^{-\lambda} e^{\lambda e^s} \\
        &= e^{\lambda (e^s - 1)}
\end{align*}
Now, if $X$ and $Y$ are independent Poisson R.V. with parameters
$\lambda_1$ and $\lambda_2$, respectively, the MGF of their sum is given by
\begin{align*}
    \phi_{X+Y}(s) &= \phi_X(s) \phi_Y(s) \\
        &= e^{\lambda_1(e^s - 1)} e^{\lambda_2 (e^s - 1)} \\
        &= e^{\lambda_1 (e^s-1) + \lambda_2(e^s - 1)} \\
        &= e^{(\lambda_1 + \lambda_2)(e^s - 1)}
\end{align*}

Since two distributions are identical if and only if their MGFs are equivalent,
it follows that the sum $X+Y$ is a Poisson distribution with parameter $\lambda_1 + \lambda_2$.


\section{Problem 6}
Since $X$ and $Y$ are independent, their joint PDF, $f_{X,Y}$ is the product of their
marginal PDFs, $f_X$ and $f_Y$:
\begin{equation*}
    f_{X,Y}(x, y) = f_X(x) f_Y(y).
\end{equation*}
We find the CDF of $Z$ by integrating over the joint PDF:
\begin{align*}
    F_Z(z) = {\bf P}(Z \leq z) &= {\bf P}(Y \leq z X) \\
        &= \int_{-\infty}^{\infty} \biggl(
            \int_{-\infty}^{zx} f_{X,Y}(x, y) dy
        \biggr) dx \\
        &= \biggl\{
            \begin{array}{ll}
                0 &\quad \text{if } z < 0 \\
                \int_{0}^{\infty} f_X(x) F_Y(xz) dx & \quad \text{otherwise}
            \end{array}
        \biggr.,
\end{align*}
where $F_Y$ is the marginal CDF of $Y$. Then, simplifying the case where $z \geq 0$, we have:
\begin{align*}
    F_Z(z) &= \int_{0}^{\infty} \lambda_1 e^{-\lambda_1 x} \biggl(
        1 - e^{-\lambda_2zx}
    \biggr) dx \\
    &= \biggl(
        -e^{-\lambda x} + \frac{\lambda_1}{\lambda_1 + \lambda_2 z} e^{-(\lambda_1 + \lambda_2 z)x}
    \biggr) \eval_{x=0}^\infty \\
    &= 1 - \frac{\lambda_1}{\lambda_1 + \lambda_2 z} \\
    &= \frac{\lambda_2 z}{\lambda_1 + \lambda_2 z}.
\end{align*} 

Using the quotient rule to differentiate yields the PDF
\begin{align*}
    f_Z(z) = \frac{dF_Z(z)}{dz} &= \frac{
        (\lambda_1 + \lambda_2z) \lambda_2 - (\lambda_2 z) \lambda_2
    }{(\lambda_1 + \lambda_2z)^2} \\
    &= \frac{\lambda_1 \lambda_2}{(\lambda_1 + \lambda_2 z)^2}
\end{align*}
which holds when $z\geq0$. For $z<0$, $f_Z(z) = 0$.

\section{Problem 7}

\begin{letterenum}
\item
Since on day 0 the weather is $A$, we have
\begin{align*}
    P_0(A) &= 1 \\
    P_0(B) = P_0(C) &= 0 \\
\end{align*}

We express $P_n(B)$ recursively in terms of the weather on day $n-1$. Let
$p=0.9$ be the probability that the weather alternates between $A$ and $B$. Since
only weather $A$ on day $n-1$ can possibly result in weather $B$ on day $n$, the
total probability law implies:
\begin{equation*}
    P_n(B) = {\bf P}(A \text{ transitions to } B | A \text{ on day } n - 1)
        {\bf P}(A \text{ on day } n-1) = p P_{n-1}(A)
\end{equation*}
and by symmetry we have
\begin{equation*}
    P_n(A) = p P_{n-1}(B).
\end{equation*}
We enumerate the first few days
\begin{align*}
    P_1(A) = 0 &\qquad P_1(B) = p \\
    P_2(A) = p^2 &\qquad P_2(B) = 0 \\
\end{align*}
and, by inspection, see that
\begin{align*}
    P_n(A) &= \biggl\{
        \begin{array}{ll}
            0 &\quad \text{if }n\text{ is odd} \\
            p^n &\quad \text{otherwise}
        \end{array}
    \biggr. \\
    P_n(B) &= \biggl\{
        \begin{array}{ll}
            p^n &\quad \text{if }n\text{ is odd} \\
            0 &\quad \text{otherwise}
        \end{array}
    \biggr.
\end{align*}
A formal proof by induction is straightforward and omitted.

\item The events that the weather is $A$, $B$, and $C$ on day $n$ form a 
partition of the sample space so that
\begin{align*}
    P_n(C^c) &= P_n(A \cup B) \\
        &= P_n(A) + P_n(B) \\
        &= \biggl\{
            \begin{array}{ll}
                p^n + 0 &\quad \text{if }n \text{ is even} \\
                0 + p^n &\quad \text{if }n \text{ is odd} \\
            \end{array}
        \biggr. \\
        &= p^n
\end{align*}
where $P_n(C^c)$ denotes the complement of weather $C$ on day $n$ (the probability that
the weather is not of type $C$).

\item Since $0<p<1$, the probability of each weather in the limit $n\to\infty$ is
clear from the previous formulas
\begin{align*}
    \lim_{n\to\infty} P_n(A) = \lim_{n\to\infty} P_n(B) &= 0 \\
    \lim_{n\to\infty} P_n(C) = \lim_{n\to\infty} (1 - P_n(C^c)) &= 1
\end{align*}

\item On each day there is probability $q=1 - p = 0.1$ that the weather transitions to
type $C$. The number of days until this occurs the first time is therefore a geometric
random variable $Z$ with parameter $q$ and its expected value is ${\bf E}[Z] = 1/q=10$ days.

\end{letterenum}

\section{Problem 8}

Let $A$ be the event that $u < 1/3$. By the total probability law, we can express
the PDF $f_X$ in terms of the conditional PDFs $f_{X|A}$ and $f_{X|A^c}$:
\begin{equation*}
    f_X(x) = f_{X | A}(x) {\bf P}(A) + f_{X | A^c}(x) {\bf P}(A^c)
\end{equation*}

Given $A$, the conditional PDF is exponential with parameter $\lambda$:
\begin{equation*}
    f_{X | A}(x) = \biggl\{
        \begin{array}{ll}
            \lambda e^{-\lambda x} &\quad \text{if }x \geq 0 \\
            0 &\quad \text{otherwise.}
        \end{array}
    \biggr.
\end{equation*}
Given $A^c$, the conditional PDF is uniformly distributed between 1/3 and 1:
\begin{equation*}
    f_{X|A^c}(x) = \biggl\{
        \begin{array}{ll}
            \frac{3}{2} &\quad \text{if } \frac{1}{3} \leq x < 1 \\
            0 &\quad \text{otherwise.}
        \end{array}
    \biggr.
\end{equation*}
Since the probability of $A$ is ${\bf P}(A) = 1/3$, our PDF is therefore given
by:
\begin{equation*}
    f_X(x) = \biggl\{
        \begin{array}{ll}
            \frac{\lambda}{3} e^{-\lambda x} + 1 &\quad \text{if }\frac{1}{3} \leq x < 1 \\
            \frac{\lambda}{3} e^{-\lambda x} &\quad \text{if }0 \leq x < \frac{1}{3} \text{ or } x \geq 1 \\
            0 &\quad \text{otherwise.}
        \end{array}
    \biggr.
\end{equation*}

\section{Problem 9}

\begin{letterenum}
\item To calculate the probability that $B>A^2C$, we integrate the joint PDF over
the volume where this inequality holds true. Since $A$, $B$, and $C$ are independent,
the joint PDF $f_{A,B,C}$ is given by:
\begin{align*}
    f_{A,B,C}(a, b, c) &= f_A(a) f_B(b) f_C(c) \\
        &= \biggl\{
            \begin{array}{ll}
                1 & \quad \text{if } 0\leq a, b, c < 1 \\
                0 & \quad \text{otherwise}
            \end{array}
        \biggr.
\end{align*}
The desired probability is computed as follows:
\begin{align*}
    P(B > A^2C) &= \int_{0}^{1} \int_{0}^{1} \biggl(
        \int_{a^2 c}^{1} 1 db
    \biggr) da dc \\
        &= \int_{0}^{1} \int_0^1 (1 - a^2c) da dc \\
        &= \int_{0}^{1} \biggl(
            a - \frac{a^3}{3} c
        \biggr)\eval_{a=0}^1 dc \\
        &= \int_{0}^{1} 1 - \frac{1}{3}c dc \\
        &= c - \frac{1}{6} c^2 \eval_0^1 \\
        &= \frac{5}{6}
\end{align*}

\item We run the following code in a Python console to simulate the same event:
\begin{lstlisting}
Python 3.8.10 (default, Nov 22 2023, 10:22:35)
[GCC 9.4.0] on linux
Type "help", "copyright", "credits" or "license" for more information.
>>> import random
>>> def trial():
...     a = random.uniform(0, 1)
...     b = random.uniform(0, 1)
...     c = random.uniform(0, 1)
...     return b > (a**2 * c)
...
>>> n = 100000
>>> results = [trial() for _ in range(n)]
>>> sum(results) / n
0.83259

>>> error = sum(results) / n - 5/6
>>> print(f"Error: {error:.5f}")
Error: -0.00074
\end{lstlisting}

Our simulation appears to converge to the same result as in part a).
\end{letterenum}





\end{document}
